\documentclass[11pt]{article}
\usepackage[a4paper,margin=25mm]{geometry}
\usepackage{amsmath,amssymb,amsthm,mathtools,microtype,graphicx,needspace}
\usepackage[colorlinks=true,linkcolor=blue,citecolor=blue,urlcolor=blue]{hyperref}
\newtheorem{theorem}{Theorem}
\newtheorem{proposition}[theorem]{Proposition}
\newtheorem{lemma}[theorem]{Lemma}
\newtheorem{corollary}[theorem]{Corollary}
\theoremstyle{remark}\newtheorem{remark}[theorem]{Remark}
\newcommand{\D}{\mathbb D}
\newcommand{\T}{\mathbb T}
\newcommand{\R}{\mathbb R}
\newcommand{\C}{\mathbb C}
\newcommand{\Rea}{\operatorname{Re}}
\newcommand{\Ima}{\operatorname{Im}}
\newcommand{\avg}[1]{\left\langle #1\right\rangle}
\newcommand{\ee}{\mathrm e}
\title{A sharp polynomial Lyapunov threshold\\for a stable cubic family}
\author{Lluis Eriksson}
\date{18 October 2026}
\begin{document}
\maketitle
\begin{abstract}
We determine exactly when the globally asymptotically stable cubic
family introduced by Liu and Fitzsimmons admits polynomial Lyapunov
certificates. In their normalization, with damping \(a>0\) and angular
forcing \(b\in\R\), the criterion is
\[
 |b|^{2/3}<1+a^{2/3}.
\]
The same strict inequality characterizes positive homogeneous
polynomial certificates with weak nonpositive Lie derivative, strict
polynomial certificates, and local real-analytic certificates.
Equality admits none of them. Below the threshold we construct a
certificate as a finite product of positive quadratic forms.
The proof pairs positive harmonic functions with an explicit analytic
square-root test and constructs a complementary limiting harmonic
pair at the boundary. Poisson smoothing and equal-weight quantiles
then produce a polynomial with a strict margin.
At a fixed nonzero boundary point, an explicit sufficient degree is
\(O_b(\delta^{-3})\), where \(\delta\) is the damping excess, while
the minimum degree diverges; no matching degree rate is asserted.
The threshold persists for every single even angular harmonic and
nonzero angular speed. Classical binary-form factorization gives
the corresponding sum-of-squares consequence without increasing the
Lyapunov degree. The family, its global stability and the original
counterexample to a polynomial converse theorem are prior results.
Our classification concerns certificates, not a change in dynamical
stability.
\end{abstract}

\section{The family and the exact certificate region}
For \(a>0\) and \(b\in\R\), consider
\begin{equation}\label{eq:field}
 \begin{aligned}
 \dot x&=-a(x^2+y^2)x+b(x^2-y^2)x-(x^2+y^2)y,\\
 \dot y&=-a(x^2+y^2)y+b(x^2-y^2)y+(x^2+y^2)x .
 \end{aligned}
\end{equation}
In polar coordinates \(x=r\cos\theta,y=r\sin\theta\), this is
\begin{equation}\label{eq:polar}
 \dot r=r^3(-a+b\cos2\theta),\qquad \dot\theta=r^2 .
\end{equation}
Liu and Fitzsimmons introduced this family and established its
global asymptotic stability, together with nonexistence of polynomial
and local analytic certificates in the region
\(|b|\ge2(a+1)\) \cite{liu}. Their work already refutes the homogeneous
polynomial Lyapunov converse. We determine the exact certificate
region of their family.

For a vector field \(f\), write \(L_fV=\nabla V\cdot f\).
A \emph{weak certificate} means \(V(0)=0\), \(V>0\) away from the
origin, and \(L_fV\le0\). A \emph{strict certificate} has
\(L_fV<0\) away from the origin. For a local analytic certificate
these conditions are required on a punctured neighborhood of zero.
All homogeneous certificates below have positive degree.

\Needspace{15\baselineskip}
\begin{theorem}\label{thm:main}
For \eqref{eq:field}, the following are equivalent:
\begin{enumerate}
\item \( |b|<(1+a^{2/3})^{3/2}\).
\item There exists a positive definite homogeneous polynomial weak
certificate.
\item There exists a strict homogeneous polynomial certificate that
is a finite product of positive definite real quadratic forms.
\item There exists a local real-analytic weak certificate.
\end{enumerate}
Thus equality at the threshold is excluded even for weak certificates.
Every system in \eqref{eq:field} is nevertheless globally
asymptotically stable.
\end{theorem}

The stability assertion is recalled from \cite{liu}, not new.
For orientation, its nonpolynomial certificate is
\[
 H(r,\theta)=r^2\exp(-b\sin2\theta),\qquad
 L_fH=-2ar^2H.
\]
It is positive, coercive, continuously differentiable at the origin,
and smooth away from it. The new assertion in Theorem~\ref{thm:main}
is the exact algebraic and analytic threshold.

\begin{figure}[t]
\centering
\includegraphics[width=.78\textwidth]{images/phase-boundary.pdf}
\caption{Certificate regions for the stable cubic family. All displayed
parameters with \(a>0\) give globally asymptotically stable dynamics.
The solid curve is excluded from the polynomial and local analytic
certificate region. The dashed line marks the sufficient
nonexistence condition in \cite{liu}; it is not a stability boundary.}
\label{fig:phase}
\end{figure}

Unbounded required polynomial degree in fixed dimension and field
degree was already established by Ahmadi and Parrilo
\cite[Theorem 4.3]{ahmadi-parrilo}. Ahmadi and El Khadir
\cite[Proposition 6.1]{ahmadi-khadir} also compare polynomial and
rational certificate complexity. Our degree result concerns approach
to the explicit boundary above. Compactness, positive harmonic
representations, and the binary-form sum-of-squares principle are
classical tools, not new general converse theorems.

\section{Polynomial certificates as positive harmonic data}
Use normalized circle averages
\(\avg v=(2\pi)^{-1}\int_0^{2\pi}v(\phi)\,d\phi\).
The circular Hilbert transform \(\mathcal H\) is normalized by
\[
 \mathcal H\cos(k\phi)=\sin(k\phi),\quad
 \mathcal H\sin(k\phi)=-\cos(k\phi),\quad
 \mathcal H1=0,\qquad k\ge1.
\]
On smooth real functions it is skew-adjoint for this average.
The formulas used below can equally be verified termwise from
Fourier series.

\begin{lemma}\label{lem:harmonic}
Let \(P\) be a positive definite real homogeneous polynomial of
degree \(2N\), \(N\ge1\), and set
\(p(\theta)=P(\cos\theta,\sin\theta)\).
There exist \(\alpha_1,\ldots,\alpha_N\in\D\) such that
\[
 h(z)=\frac1N\sum_{j=1}^N\frac{1+\alpha_jz}{1-\alpha_jz},
 \qquad
 \frac{p'(\theta)}{2Np(\theta)}
       =\Ima h(\ee^{2i\theta}).
\]
Here \(h(0)=1\), and \(u=\Rea h\) has a strictly positive minimum
on the closed disk. The weak certificate condition is equivalent to
\begin{equation}\label{eq:weak}
 w(\phi):=a-b\cos\phi-\mathcal Hu(\phi)\ge0.
\end{equation}
\end{lemma}
\begin{proof}
The strict Fej\'er--Riesz factorization, in the form used in
\cite[Lemmas III.2--III.3]{liu}, writes
\[
 p(\theta)=|C|^2\prod_{j=1}^L
        |1-\alpha_j\ee^{2i\theta}|^2,
 \qquad L\le N,\quad |\alpha_j|<1.
\]
Pad the list with zero parameters to length \(N\).
Each kernel has real part
\[
 \Rea\frac{1+\alpha z}{1-\alpha z}
       =\frac{1-|\alpha z|^2}{|1-\alpha z|^2}>0
       \qquad(|z|\le1).
\]
Differentiating the logarithm of a factor gives twice the imaginary
part of that kernel, proving the angular identity.
The imaginary part of \(h\) is \(\mathcal Hu\) because \(h(0)=1\).
Finally,
\[
 L_fP=2N r^{2N+2}p(\theta)
       \left[-a+b\cos2\theta+
              \frac{p'(\theta)}{2Np(\theta)}\right],
\]
which proves \eqref{eq:weak}.
\end{proof}

Conversely, an equal-weight average of these kernels yields a
polynomial of degree \(2N\). Explicitly, put
\begin{equation}\label{eq:quadratic}
 Q_\alpha(x,y)=(1+|\alpha|^2)(x^2+y^2)
       -2\Rea\alpha(x^2-y^2)+4\Ima\alpha\,xy.
\end{equation}
Its eigenvalues are \((1-|\alpha|)^2,(1+|\alpha|)^2\).
The product \(P=\prod_jQ_{\alpha_j}\) is positive definite and
satisfies the angular identity in Lemma~\ref{lem:harmonic}.
A strict margin in \eqref{eq:weak} therefore gives a strict
polynomial certificate.

\section{The strict obstruction}
Rotation by \(\pi/2\) changes the sign of \(b\), so assume \(b\ge0\).
For \(c>0\), define
\begin{equation}\label{eq:test}
 \begin{gathered}
 s=\sqrt{1+c^2},\qquad \eta=\frac{c+i}{s},\\
 M_c(z)=c^2-1+2ic\frac{1+\eta z}{1-\eta z},
 \qquad f_c(z)=\sqrt{M_c(z)}.
 \end{gathered}
\end{equation}
The Cayley transform has positive real part on \(\D\); hence
\(M_c\) takes values in the upper half-plane. We choose the
analytic square root in the first quadrant. In particular,
\begin{equation}\label{eq:fcoeff}
 f_c(0)=c+i,\qquad f_c'(0)=\frac{2ic}{s}.
\end{equation}

For \(0<\sigma<1\), write
\(f_c(\sigma\ee^{i\phi})=F_\sigma(\phi)+iG_\sigma(\phi)\).
Both functions are positive and smooth. Their means and first
coefficient give
\begin{equation}\label{eq:testmeans}
 \avg{F_\sigma}=c,\quad \avg{G_\sigma}=1,\quad
 \avg{G_\sigma\cos\phi}=\frac{\sigma c}{s},\quad
 F_\sigma+\mathcal HG_\sigma=c.
\end{equation}
If \(u,w\) come from a weak polynomial certificate, then
\(\avg u=1\), \(u\ge\varepsilon>0\), and \(w\ge0\).
Using \eqref{eq:weak} and skew-adjointness,
\begin{align}
 \avg{uF_\sigma+wG_\sigma}
 &=\avg{u(F_\sigma+\mathcal HG_\sigma)}
      +a\avg{G_\sigma}-b\avg{G_\sigma\cos\phi}\nonumber\\
 &=c+a-\frac{b\sigma c}{s}\ge\varepsilon c.
 \label{eq:pairing}
\end{align}
Every function in this computation is smooth for \(\sigma<1\);
no singular boundary pairing is assumed.
Letting \(\sigma\uparrow1\) yields
\begin{equation}\label{eq:dual}
 b\le\frac{(a+c)s}{c}-\varepsilon s
      <\frac{(a+c)\sqrt{1+c^2}}c.
\end{equation}
The derivative of the last expression is
\[
 \frac{c^3-a}{c^2\sqrt{1+c^2}}.
\]
Its minimum occurs at \(c=a^{1/3}\), with value
\((1+a^{2/3})^{3/2}\). This proves the strict necessary
inequality, including nonexistence at equality. The positive
minimum of \(u\), rather than mere nonnegativity, is what
excludes the boundary.

\section{An explicit limiting pair on the boundary}
The converse needs more than a necessary moment inequality.
We construct positive harmonic data at the critical parameters,
then approximate them while retaining a positive margin.

Fix \(c>0\), with \(s,\eta,f_c\) as above, and put
\[
 a_0=c^3,\qquad b_0=s^3.
\]
Define
\begin{align}
 q_c(z)&=1-c^4+\frac{b_0}{2}
          \bigl[(c-i)z+(c+i)z^{-1}\bigr],\label{eq:q}\\
 h_c(z)&=ia_0-\frac{ib_0}{2}(z+z^{-1})
                    +i\frac{q_c(z)}{f_c(z)}.\label{eq:h}
\end{align}

\begin{lemma}\label{lem:critical}
The function \(h_c\) extends analytically across zero, continuously
to \(\overline\D\), and satisfies \(h_c(0)=1\).
Its boundary data obey
\[
 u_c=\Rea h_c\ge0,\qquad
 w_c=a_0-b_0\cos\phi-\Ima h_c(\ee^{i\phi})\ge0.
\]
Their positivity sets are disjoint arcs.
\end{lemma}
\begin{proof}
The coefficient of \(z^{-1}\) in \eqref{eq:h} cancels because
\(f_c(0)=c+i\). Its constant coefficient is
\[
 ia_0+\frac{i(1-c^4)}{f_c(0)}
       -\frac{ib_0}{2}\frac{f_c'(0)}{f_c(0)}=1,
\]
by \eqref{eq:fcoeff}. No other interior singularity occurs,
since \(f_c\) has no zero in \(\D\).

Two exact rational identities determine the boundary behavior:
\begin{align}
 M_c(z)&=s^2\frac{\eta^2-\bar\eta z}{1-\eta z},
       \label{eq:Mfactor}\\
 q_c(z)&=-\frac{s^2}{2}M_c(z)
            (1-\eta z)(1-\bar\eta/z).\label{eq:qfactor}
\end{align}
On the circle, \(M_c\) and \(q_c\) are real away from the
exceptional points, and
\[
 q_c(z)=-\frac{s^2}{2}M_c(z)|1-\eta z|^2.
\]
On the arc where \(M_c>0\), the chosen boundary square root
is real and positive; \(q_c<0\). Then \(u_c=0\) and
\(w_c=-q_c/f_c\ge0\).
On the complementary arc \(M_c<0\), write \(f_c=iG\), \(G>0\).
There \(q_c>0\), so \(u_c=q_c/G\ge0\) and \(w_c=0\).

For completeness, \(M_c\) has a simple boundary zero at
\(z=\eta^3\) and a simple boundary pole at \(z=\bar\eta\).
They are distinct for every finite \(c>0\).
The Laurent polynomial \(q_c\) vanishes at both points.
Thus \(q_c/f_c\) is respectively of orders \(1/2\) and \(3/2\),
and extends continuously with value zero.
The boundary inequalities therefore extend to both exceptional
points. The minimum principle gives \(\Rea h_c\ge0\) on
\(\overline\D\), strictly positive in the interior because
\(h_c(0)=1\).
\end{proof}

This limiting pair is not a polynomial certificate at
\((a_0,b_0)\). Its real part vanishes on an arc, unlike the
strictly positive boundary real part in Lemma~\ref{lem:harmonic}.

\section{Constructing a finite polynomial certificate}
Suppose first \(b>1\) and \(b<(1+a^{2/3})^{3/2}\).
Choose
\begin{equation}\label{eq:a0}
 c=\sqrt{b^{2/3}-1},\qquad
 a_0=(b^{2/3}-1)^{3/2}<a,\qquad b_0=b.
\end{equation}
Poisson smoothing of Lemma~\ref{lem:critical} gives, for
\(h_\sigma(z)=h_c(\sigma z)\),
\[
 \Ima h_\sigma(\ee^{i\phi})+b\sigma\cos\phi\le a_0.
\]
Consequently
\begin{equation}\label{eq:smoothmargin}
 \Ima h_\sigma(\ee^{i\phi})+b\cos\phi
       \le a_0+b(1-\sigma)<a
\end{equation}
when \(\sigma<1\) is sufficiently close to one.

For \(0\le b\le1\), instead use \(h(z)=1-ibz\).
Its boundary real part is nonnegative, \(h(0)=1\), and
\(\Ima h(\ee^{i\phi})+b\cos\phi=0\).
Replacing \(z\) by \(\sigma z\) gives a strict margin whenever
\(b(1-\sigma)<a\). This treats \(b=0\) and \(b=1\)
without using the singular endpoint \(c=0\) in the critical formula.

In both cases the unsmoothed \(h\) is analytic, continuous on the
closed disk, has nonnegative real part and value 1 at zero.
Poisson's formula and its harmonic conjugate imply
\begin{equation}\label{eq:herglotz}
 h(\sigma z)=\int_0^{2\pi}
       \frac{1+\sigma\ee^{-it}z}{1-\sigma\ee^{-it}z}\,d\mu(t),
 \qquad
 d\mu(t)=\Rea h(\ee^{it})\,\frac{dt}{2\pi}.
\end{equation}
The measure is a probability measure. The imaginary constant is
zero because \(h(0)=1\).
For fixed \(\sigma<1\) the kernel is continuous jointly in \(t\)
and \(|z|\le1\). Approximate \(\mu\) by equal-weight empirical
measures; uniform continuity makes the resulting approximation
uniform on the closed disk. Thus, for sufficiently large \(N\),
\[
 h_N(z)=\frac1N\sum_{j=1}^N
        \frac{1+\alpha_jz}{1-\alpha_jz},
 \qquad \alpha_j=\sigma\ee^{-it_j},
\]
retains the strict margin in \eqref{eq:smoothmargin}.
The product of the quadratic forms \eqref{eq:quadratic} is
then a strict certificate. The next proposition gives an explicit
choice of a sufficient \(N\).

\begin{proposition}[A sufficient degree bound]\label{prop:degree}
Let \(b>1\), \(a_0\) be as in \eqref{eq:a0}, and
\(\delta=a-a_0>0\). Set
\[
 \tau=\min\{1/2,\delta/(4b)\},\qquad
 \sigma=1-\tau,\qquad
 N=\max\left\{1,\left\lceil
                  \frac{16\pi\sigma}{\delta\tau^2}
                    \right\rceil\right\}.
\]
There is a strict polynomial certificate of degree \(2N\)
whose normalized angular margin
\(a-b\cos2\theta-p'/(2Np)\) is at least \(\delta/2\).
For \(0<\delta\le2b\), one may use
\[
 N\le\left\lceil\frac{256\pi b^2}{\delta^3}\right\rceil.
\]
This is a sufficient bound, not the minimum degree.
\end{proposition}
\begin{proof}
The critical measure in \eqref{eq:herglotz} is nonatomic.
Partition \([0,2\pi]\) into \(N\) successive quantile intervals
\(I_j\), each of mass \(1/N\), and choose \(t_j\in I_j\).
Intervals may cross zero-density gaps; their lengths still sum
to \(2\pi\).
For the kernel \(K_\sigma(t,z)\) in \eqref{eq:herglotz},
\[
 \sup_{t,|z|\le1}|\partial_tK_\sigma(t,z)|
       \le\frac{2\sigma}{(1-\sigma)^2}.
\]
The integral-minus-empirical error is therefore at most
\[
 \sum_j\int_{I_j}\frac{2\sigma}{\tau^2}|I_j|\,d\mu
       =\frac{4\pi\sigma}{N\tau^2}\le\frac\delta4.
\]
The smoothing loss is \(b\tau\le\delta/4\), leaving margin
at least \(\delta/2\).
When \(\delta\le2b\), insert \(\tau=\delta/(4b)\) and
\(\sigma\le1\) to obtain the final bound.
\end{proof}

Locating the quantiles with rigorous numerical enclosures is an
additional computational task. Proposition~\ref{prop:degree} is an
analytic construction and degree bound, not a claim that the
floating-point experiment below implements a certified algorithm.

\section{Local analytic certificates and degree divergence}
We complete the equivalences of Theorem~\ref{thm:main}.
Let \(V\) be a local real-analytic weak certificate, and let
\(P_m\) be its first nonzero homogeneous Taylor term.
Positivity of \(V\) implies \(P_m\ge0\), and comparison of the
first Lie-derivative terms gives \(L_fP_m\le0\).
On the circle its profile satisfies
\begin{equation}\label{eq:leading}
 p_m'(\theta)+m(-a+b\cos2\theta)p_m(\theta)\le0.
\end{equation}
If \(p_m(\theta_0)=0\), multiply by the positive integrating
factor of \eqref{eq:leading}. The resulting function is
nonnegative and nonincreasing from zero over one full turn.
It must vanish throughout the turn, contradicting \(P_m\ne0\).
Hence \(P_m\) is positive definite. Its degree is even, and
the strict obstruction applies. This is the leading-term
mechanism already used in \cite{liu}.
Conversely, the constructed polynomial is analytic.
Together with the preceding sections this proves
Theorem~\ref{thm:main}.

\begin{proposition}\label{prop:blowup}
For fixed \(b>1\), the minimum degree among positive homogeneous
polynomial weak certificates tends to infinity as
\(a\downarrow a_0=(b^{2/3}-1)^{3/2}\).
\end{proposition}
\begin{proof}
Otherwise select a sequence of certificates with bounded degrees,
and a subsequence of one fixed even degree \(2N\).
Normalize their circle restrictions to maximum one.
Their Fourier coefficients are bounded, so a subsequence converges
with its derivative to a nonnegative trigonometric polynomial
\(p\) of maximum one. The weak inequality passes to the limit
at \(a=a_0\). The integrating-factor argument above forces
\(p>0\). It is the circle restriction of a homogeneous polynomial
of degree \(2N\), even if its Fourier degree drops.
This would give a weak certificate at the excluded boundary,
a contradiction.
\end{proof}

This compactness method has antecedents in
\cite[Theorem 4.3]{ahmadi-parrilo}.
Propositions~\ref{prop:degree} and \ref{prop:blowup}
do not give matching upper and lower degree rates.

\section{Single harmonics, angular speed and sum of squares}
The same extremal argument applies to a family in every odd degree.
Let \(k\ge1\), \(a>0\), \(b\in\R\), and \(\omega\ne0\), and consider
\begin{equation}\label{eq:higher}
 \dot r=r^{2k+1}(-a+b\cos2k\theta),
 \qquad \dot\theta=\omega r^{2k}.
\end{equation}
It is a polynomial vector field of degree \(2k+1\):
\[
 f(x,y)=
 \bigl[-a(x^2+y^2)^k+b\Rea(x+iy)^{2k}\bigr](x,y)
       +\omega(x^2+y^2)^k(-y,x).
\]

\begin{corollary}\label{cor:higher}
For \eqref{eq:higher}, homogeneous polynomial weak, strict
polynomial, and local real-analytic weak certificates exist
exactly when
\begin{equation}\label{eq:highercriterion}
 |b|^{2/3}<a^{2/3}+|\omega|^{2/3}.
\end{equation}
All these fields are globally asymptotically stable.
For \(|b|>|\omega|\), the degree bound of
Proposition~\ref{prop:degree}, applied to
\(a/|\omega|,|b|/|\omega|\), becomes \(2kN\).
\end{corollary}
\begin{proof}
Positive time rescaling, reflection and rotation reduce to
\(\omega=1\), \(b\ge0\).
For an arbitrary certificate of degree \(2N\),
Lemma~\ref{lem:harmonic} gives
\(w=a-b\cos(k\phi)-\mathcal Hu\ge0\).
Use the test \(f_c(\sigma z^k)\) in place of
\(f_c(\sigma z)\). Its means are still \(c,1\),
its real and imaginary parts are positive, and
\[
 \avg{G_\sigma\cos(k\phi)}=\frac{\sigma c}{s},
 \qquad F_\sigma+\mathcal HG_\sigma=c.
\]
The pairing \eqref{eq:pairing} and its strict conclusion are
unchanged. This necessity does not assume that the candidate
polynomial has \(k\)-fold angular symmetry.

For sufficiency take \(L\) nodes from the cubic construction and set
\[
 P(r,\theta)=r^{2kL}
       \prod_{j=1}^L|1-\alpha_j\ee^{2ik\theta}|^2.
\]
Each factor is the positive binary form
\[
 |(x-iy)^k-\alpha_j(x+iy)^k|^2
 =(1+|\alpha_j|^2)(x^2+y^2)^k
       -2\Rea\bigl[\alpha_j(x+iy)^{2k}\bigr].
\]
The product has degree \(2kL\) and satisfies
\[
 \frac{p'(\theta)}{2kL p(\theta)}
       =\Ima h_L(\ee^{2ik\theta}).
\]
Thus the strict margin persists.
The local analytic argument is unchanged because the angular
speed does not vanish away from zero.
Finally,
\[
 H=r^2\exp\left[-\frac{b}{k\omega}\sin2k\theta\right],
 \qquad L_fH=-2ar^{2k}H.
\]
Its positive bounded angular factor gives coercivity and
\(C^1\) regularity at zero. Decrease of \(H\) bounds all
trajectories and prevents finite-time escape; the displayed
derivative and comparison with \(r^2\) imply global
asymptotic stability.
\end{proof}

For \(k=1\), replacing \(b\cos2\theta\) by
\(b\cos2\theta+d\sin2\theta\) merely rotates its phase.
The criterion becomes
\[
 (b^2+d^2)^{1/3}<a^{2/3}+|\omega|^{2/3}.
\]
No result for \(\omega=0\) or a general sum of different
angular harmonics follows from this argument.

\begin{corollary}[Classical binary-form SOS consequence]\label{cor:sos}
The strict certificate region in Theorem~\ref{thm:main} and
Corollary~\ref{cor:higher} is also the region admitting a
homogeneous certificate \(P\) such that \(P\) and \(-L_fP\)
are positive definite sums of squares.
The same sufficient Lyapunov degree bounds apply.
\end{corollary}
\begin{proof}
Any positive definite real binary form \(U\) of degree \(2m\)
is a sum of two squares of real forms of degree \(m\).
Indeed \(U(x,1)\) has positive leading coefficient and no real
roots. Splitting its roots into nonreal conjugate pairs gives
\(U(x,1)=F(x)\overline{F(x)}\). Homogenizing \(F\) yields
\(U=(\Rea F)^2+(\Ima F)^2\).
Apply this classical fact separately to the strict constructed
\(P\) and \(-L_fP\). The reverse implication is immediate.
\end{proof}
This binary SOS principle and broader converse results are
discussed in \cite[Section 5 and Theorem 5.3]{ahmadi-parrilo};
no new general SOS theorem is claimed.
If the vector field has rational coefficients, a sufficiently
small rational perturbation of the coefficient vector of \(P\)
preserves both strict inequalities on the compact unit circle.
Thus a rational-coefficient polynomial certificate exists as well.
This does not assert a two-square decomposition over the rationals.

\section{Verification and limitations}
The proof uses explicit analytic identities and inequalities, not
a numerical optimization result. Supplementary author computations
check the rational factorization of the critical pair, the value
\(h_c(0)=1\), and five boundary parameter instances.
One finite product of 2048 quadratic factors was evaluated at
8192 angular samples for \(a=1.2\), \(b=2\sqrt2\),
\(\sigma=0.98\). The sampled normalized margin was about
\(0.1545\). This is a floating-point diagnostic candidate,
not a certified global positivity check or a minimum-degree result.

A separate diagnostic solves 18 necessary finite Toeplitz
relaxations at three damping values and six truncation orders.
The stored field named \texttt{degree} in that exploratory program
means moment truncation order, not Lyapunov polynomial degree.
Nine solver results are marked \texttt{optimal\_inaccurate};
these labels are preserved. No solver status is used as a proof
of feasibility, nonexistence or sharpness.

The same involved model reviewer separately checked the proof
dossier and additional symbolic identities, finite Hilbert
pairings and high-precision branch-point profiles.
This is not blind or human peer review, an independent additional
vote, or formal verification. No Lean build of the earlier
source or this argument was performed.
The quantitative construction is analytic; a rigorously enclosed
implementation of its quantile selection is not supplied.

The comparison uses the complete mathematical text of
Liu--Fitzsimmons v1 and the cited degree-complexity and SOS portions
of the earlier works, not complete readings of the latter.
Targeted searches did not locate this boundary; worldwide priority
is unconfirmed, particularly for the associated harmonic extremal
problem. Minimum-degree asymptotics and general multiharmonic
classification remain open here. AI tools assisted this work.

\small
\begin{thebibliography}{9}
\bibitem{liu}
J. Liu and M. Fitzsimmons,
\emph{A globally asymptotically stable planar homogeneous polynomial
vector field with no polynomial Lyapunov function},
arXiv:2607.16171v1, 17 July 2026.
\url{https://arxiv.org/abs/2607.16171v1}.
\bibitem{ahmadi-parrilo}
A. A. Ahmadi and P. A. Parrilo,
\emph{Stability of polynomial differential equations:
complexity and converse Lyapunov questions}, arXiv:1308.6833.
Theorem numbering refers to the author-hosted manuscript:
\url{https://web.mit.edu/~a_a_a/Public/Publications/Selected_Papers/dTAC.pdf}.
\bibitem{ahmadi-khadir}
A. A. Ahmadi and B. El Khadir,
\emph{On algebraic proofs of stability for homogeneous vector fields},
IEEE Trans. Automat. Control \textbf{65} (2020), 325--332.
The comparison uses Proposition 6.1 and Lemma 6.2 of the 2018 preprint:
\url{https://optimization-online.org/wp-content/uploads/2018/03/6502.pdf}.
\end{thebibliography}
\end{document}
